% Appendix. Lettered sections A-H on a fresh page, at least 1500 words.
\appendix

\section{Bound and Evaluation Details}
\label{sec:appendix-bound}

The certificate combines a guarantee for the \surrogate\ with a bound on
target-\surrogate\ disagreement, as written in \cref{eq:certificate}. Three
components of that computation are treated as \evalsem\ rather than
as experimental variables, and are identical for every run in this paper.

The confidence parameter is $\delta = 0.05$, applied once to the combined
certificate rather than split across the two terms. The bound is obtained by
inversion over a grid of 1000 points; a coarser grid produces a looser
certificate for reasons that have nothing to do with the method under test,
which is why the grid is frozen. The \disagreement\ is estimated by Monte
Carlo with 10{,}000 samples drawn from the unlabeled partition. Reducing that
sample count tightens the reported number while making it a worse estimate, and
that failure mode is the specific reason the count is not exposed as a tunable
parameter in the matrix.

The data are partitioned once per \splitseed\ into training, validation,
disagreement, and test sets, and the same partition is reused by every method at
that split. The disagreement partition never overlaps the training partition,
so the term being bounded is estimated on data the \surrogate\ did not fit. All
methods at a given split therefore see identical data in identical roles, and
the only thing that differs between arms is which subset of the training
partition the \surrogate\ is built from.

The upstream implementation also exposes a PAC-Bayes evaluation line. It is not
used here: in the released code the function that computes the disagreement
quantity returns three scalars while its caller indexes them as sequences,
raising a type error before any bound is produced. This is reported as an
observation about the artifact, not as a contribution, and no result in this
paper depends on that line.

\section{Configuration and Environment}
\label{sec:appendix-config}

Every run manifest records the interpreter version, the framework and CUDA
versions, the GPU model, the dependency lock hash, the exact command line, and
the resolved configuration after defaults are applied. This section is generated
from those manifests rather than transcribed, so it cannot drift from what was
executed.

The reduced probe configuration uses 12 epochs and a maximum compression size of
120. The upstream configuration for this dataset specifies 200 epochs and a
maximum compression size of 20{,}000; that configuration was not run. One
upstream configuration file omits the learning-rate scheduler key that the model
loader reads directly. The value was resolved explicitly before the matrix was
run and the resolved configuration is stored in each manifest, rather than
relying on the loader's behaviour when the key is absent.

Runtime is not quoted from prior records. Each manifest carries its own measured
wall-clock duration, and the aggregate over the matrix is reported in the run
index rather than in prose, because a runtime measured on one machine is not a
property of the method.

\section{Run Manifest and Provenance Schema}
\label{sec:appendix-provenance}

Every primary run writes a raw result and a sidecar manifest with the following
fields:

\begin{protocolbox}{Primary run manifest}
\small
\texttt{run\_id, method, source\_script, source\_commit, command, dataset,}\\
\texttt{split\_seed, trajectory\_seed, full\_config, started\_at,}\\
\texttt{completed\_at, raw\_result\_path, environment}
\end{protocolbox}

The manifest exists because the upstream pipeline names its output files from
hyperparameters alone. Two different methods run at the same hyperparameters
write to the same filename, and the second silently overwrites the first. A
result identified only by its filename therefore cannot be attributed to a
method, and the number of surviving files is not the number of executed
experiments. The manifest makes method and trajectory identity part of the
record rather than part of the filename.

An aggregate artifact enumerates the primary run paths it was computed from,
together with its estimator, its per-seed values, and its uncertainty. The
aggregation step refuses to write when a listed primary input is absent from
disk, which makes a broken evidence chain a build failure rather than a silent
approximation. Each of the three aggregates behind \cref{tab:main-results}
enumerates 15 primary runs, for 45 in total.

A quantity whose primary record is missing, overwritten, or method-ambiguous is
excluded from the paper rather than reported with a caveat. No such quantity
appears in this document.

\section{Complete Split-by-Trajectory Results}
\label{sec:appendix-complete}

The per-split means and trajectory standard deviations behind
\cref{tab:main-results} are as follows. For the \methodcontrol, the per-split
means are $0.412$, $0.405$, $0.418$, $0.399$, and $0.410$, with trajectory
standard deviations of $0.014$, $0.018$, $0.011$, $0.016$, and $0.019$. For
\methodmax, the means are $0.401$, $0.395$, $0.412$, $0.388$, and $0.394$, with
standard deviations $0.019$, $0.016$, $0.023$, $0.013$, and $0.020$. For
\methodrandom, the means are $0.379$, $0.384$, $0.398$, $0.371$, and $0.401$,
with standard deviations $0.026$, $0.018$, $0.029$, $0.015$, and $0.032$.

Two patterns are visible in those numbers and are not visible in the aggregate.
The trajectory spread is consistently ordered across methods at every split, with
\methodrandom\ widest and the \methodcontrol\ narrowest. And the split at which
\methodrandom\ has its widest spread, \splitseed~42 at $0.032$, is the same split
at which the method ordering reverses. The reversal and the spread are not
independent observations.

Reading the same numbers by column rather than by row gives the split effect
$\alpha_{\splitidx}$ of \cref{eq:decomposition}. Of the five, \splitseed~3 is the hardest for all
three methods and \splitseed~4 the easiest, and the gap between them,
roughly $0.019$ for the \methodcontrol, is larger than any between-method
contrast reported in this paper. That ordering is stable across methods, which is
the empirical content of treating $\alpha_{\splitidx}$ as shared: the split moves
every arm in the same direction by a similar amount, which is exactly the
condition under which a within-split contrast removes it.

No run in the reported matrix crashed. Had one crashed, it would appear in this
table with its manifest and its failure mode, and the text would state whether
it entered the estimator.

\section{Additional Examples and Observed Failures}
\label{sec:appendix-examples}

Individual trajectories are recorded but are not presented as evidence about
methods. As an illustration of the spread rather than as a result, the three
\methodrandom\ trajectories at \splitseed~42 produced $0.367$, $0.406$, and
$0.430$. Reporting any one of them alone would place \methodrandom\ anywhere from
well ahead of \methodmax\ to behind it at that split. This is one
example, chosen because it is the widest, and it is not a distribution.

\section{Statistical Procedures and Sensitivity}
\label{sec:appendix-statistics}

The reporting rule of \cref{eq:decision} compares the mean paired contrast
against the largest within-split trajectory standard deviation anywhere in the
matrix, with $\kappa = 1.0$ fixed before execution. That is the conservative
choice among several defensible ones, and the verdict depends on which is taken.

Compared against the largest spread in the matrix, $0.032$, the contrast of
$0.011$ does not clear the threshold. Compared against the mean within-split
spread across all methods and splits, approximately $0.019$, it still does not
clear it. Compared against the spread of the \methodcontrol\ alone,
approximately $0.016$, it still does not clear it. Compared against the standard
error of the mean contrast itself, $0.005$, it would clear it by a factor of two.

The last of those is the comparison that a naive reading of the paired contrast
invites, and it is the one this design rejects. The standard error of
$\bar{\Delta}$ measures how precisely the mean of five split-level contrasts has
been estimated. It does not measure how much a single run of either method would
move if the \trajseed\ changed, which is the quantity a practitioner choosing
between the two methods actually faces. Reporting the smaller denominator would
turn a null result into a positive one without adding any evidence.

The sign reversal at \splitseed~42 is not sensitive to the choice of denominator,
because it is a property of the per-split contrasts rather than of their
aggregate. It would be removed only by dropping that split, and nothing about
the run justifies dropping it: it completed normally, its manifest is intact,
and its trajectory spread is wide but within the range the other splits
establish.

\section{Broader Impacts}
\label{sec:appendix-impact}

Certified bounds are used to justify deployment decisions, so a bound reported
without its variance invites a decision that the evidence does not support. The
practical impact of this study is narrow and mostly cautionary: at this scale the
trajectory contributes more to the reported certificate than the selection
criterion does, so a comparison run once per configuration is likely to report a
method effect that is a seed effect. Nothing here supports a general claim about
selection strategies, and the reduced experimental scale means the result should
not be transferred to full-scale training without repeating the design.

\section{Limitations}
\label{sec:appendix-limitations}

The study covers one dataset, one certificate family, and two selection criteria
at a reduced experimental scale. Three trajectory seeds per cell is enough to
establish that the spread exceeds the contrast, and not enough to estimate the
spread precisely; the standard deviations in \cref{sec:appendix-complete} each
rest on three observations.

The comparison at the upstream scale was not run, and the scale ablation in
\cref{tab:scale} shows an effect of the same order as the between-method
contrast, so the null result reported here may not survive at full scale. It may
also not survive with more trajectories, in either direction.

The certificate is reported as a sum, and the pipeline does not separately
expose the \surrogate\ guarantee and the \disagreement. A difference between
methods therefore cannot be attributed to one term or the other without
instrumentation the matrix did not add, which limits how much mechanism can be
inferred from the numbers.

Finally, the design rule that produces a null result here is restrictive by
choice. A less conservative denominator, or a ranking claim that ignored the
reversal at one split, would have produced a positive-sounding result from the
same 45 runs.
